\documentclass[10pt,a4paper]{article}
\usepackage[margin=25mm]{geometry}
\usepackage{amsmath,amssymb,amsthm,booktabs,array}
\usepackage[T1]{fontenc}
\usepackage{lmodern}
\usepackage{xurl}
\usepackage[hidelinks]{hyperref}
\usepackage{microtype}
\setlength{\parskip}{4pt}
\setlength{\parindent}{1em}
\setlength{\emergencystretch}{2em}
\renewcommand{\arraystretch}{1.13}
\newtheorem{theorem}{Theorem}
\newtheorem{lemma}[theorem]{Lemma}
\newtheorem{proposition}[theorem]{Proposition}
\theoremstyle{definition}
\newtheorem{definition}[theorem]{Definition}
\theoremstyle{remark}
\newtheorem{remark}[theorem]{Remark}
\newcommand{\C}{\mathbb C}
\newcommand{\coef}[1]{[#1]P}
\title{Global identifiability of a behavioral SIR model\newline
for every integer saturation exponent $k\geq2$}
\author{Tatsuhiko Takimoto\\\small Independent researcher}
\date{Version 1.1 (draft), 22 October 2026}
\begin{document}
\maketitle
\begin{abstract}
Pant, Saucedo and Pogudin studied a behavioral SIR model in which the
transmission rate is divided by $(1+\alpha I)^k$, and conjectured that for
$k\geq2$ it is globally structurally identifiable from continuous,
noise-free observation of the number of infectious individuals $I(t)$,
with the total initial population known. We prove this for every fixed
integer $k\geq2$: generically, the parameters $\alpha,\beta,\gamma$ and the
initial values of $S$, $I$ and $R$ are uniquely determined by $I(t)$. The
proof has two parts. First, the minimal input--output equation is written
in closed form for all $k$; its coefficients determine $\alpha$, $\gamma$
and $\beta$ through rational formulas whose denominators are $k$ and
$k(k-1)\alpha^2$, so no integer $k\geq2$ is exceptional. Second, a
valuation argument shows that the two-state model obtained by eliminating
$R$ has no nonconstant rational first integral; the argument uses only
that $k\geq2$. Results of Ovchinnikov, Pogudin and Thompson then convert
the coefficient statement into global identifiability. The proof is
written in prose and is not formalized in a proof assistant. Exact
computer-algebra checks for $k=1,\dots,12$, and a separately written
check for $k=2,\dots,9$, agree with the general formulas; they are finite
checks and do not replace the proof. Non-integer exponents, an unknown
exponent, and estimation from discrete or noisy data are outside the
scope of the result.
\end{abstract}

\section{Introduction}
\label{sec:intro}
Pant, Saucedo and Pogudin examined which epidemiological quantities are
structurally identifiable in a range of compartmental models
\cite{PSP}. In their Section~3.6 they consider a behavioral extension of
the SIR model in which contacts decrease as infections increase,
\begin{equation}
 S'=-\frac{\beta SI}{(1+\alpha I)^k},\qquad
 I'=\frac{\beta SI}{(1+\alpha I)^k}-\gamma I,\qquad
 R'=\gamma I,
 \label{eq:model}
\end{equation}
observed through $y=I$, with the total initial population
$N_0=S(0)+I(0)+R(0)$ known. For $k=1$ they proved that $\alpha$,
$\beta\gamma$ and $\alpha\gamma+\beta$ are globally identifiable while
$\beta$ and $\gamma$ are only locally identifiable, the two-fold ambiguity
being $(\beta,\gamma)\mapsto(\alpha\gamma,\beta/\alpha)$
\cite[Theorem~3.8]{PSP}, \cite[Section~S8]{PSPsupp}. Using
\mbox{StructuralIdentifiability.jl}, they found global identifiability for
$k=2,3,4,5,10$ and $50$ \cite[Section~S9]{PSPsupp}, and they stated global identifiability for
$k\geq2$ as Conjecture~2 of \cite{PSP}. Their Remark~1 in
\cite[Section~S9]{PSPsupp} sketches why the property should hold for all
sufficiently large $k$: after the substitution $A=(1+\alpha I)^k/\beta$,
the parameters and initial conditions are rational functions of the
output jet, whose denominators can vanish for at most finitely many
values of $k$; the exceptional values are not determined.

The tools we use are due to Ovchinnikov, Pogudin and Thompson \cite{OPT}.
They showed that the coefficients of a monic characteristic presentation
of the input--output ideal generate the field of input--output
identifiable functions \cite[Corollary~5.7]{OPT}, and that input--output
identifiability coincides with identifiability when the model has no
rational first integrals \cite[Theorem~4.7]{OPT}. The second hypothesis
is essential here: with all three states, $S+I+R$ is a rational first
integral, so we work with the two-state model on the surface of known
$N_0$.

\paragraph{Result.}
Theorem~\ref{thm:main} proves global identifiability for every fixed
integer $k\geq2$. Its new ingredients are a closed form of the minimal
input--output equation valid for all $k$, an explicit rational inverse of
its coefficient map (Proposition~\ref{prop:coef}), and a lemma that the
reduced model has no nonconstant rational first integral
(Lemma~\ref{lem:noint}). The hypothesis $k\geq2$ is used essentially in
these two places, each time through a single, visible fact: a coefficient $k(k-1)\alpha^2\gamma$ that
vanishes for $k=1$, and a pole of order $k$ that cannot be a logarithmic
derivative when $k\geq2$. Both facts are consistent with the $k=1$
non-identifiability result of \cite{PSP}.

\section{What is proved, and how}
\label{sec:status}
Table~\ref{tab:status} lists the claims of this article and their present
basis. All proofs are written in prose; none is formalized in a proof
assistant, and the article has not undergone human expert review. Who
ran which check, and when, is recorded in the verification record that
accompanies the article (\mbox{\texttt{VERIFICATION-RECORD.md}}).
\begin{table}[ht]
\centering\small
\begin{tabular}{@{}>{\raggedright\arraybackslash}p{.30\linewidth}
>{\raggedright\arraybackslash}p{.38\linewidth}
>{\raggedright\arraybackslash}p{.24\linewidth}@{}}
\toprule
Claim & Basis & Range\\
\midrule
Input--output equation \eqref{eq:P}, its irreducibility and
characteristic presentation (Section~\ref{sec:io}) & Elementary algebra
(Gauss's lemma) and the definitions of \cite{OPT} & Every integer
$k\geq2$\\
Coefficient table and inverse map (Proposition~\ref{prop:coef}) &
Binomial theorem & Every integer $k\geq2$\\
No rational first integral (Lemma~\ref{lem:noint}) & Formal Laurent
expansion and logarithmic derivatives & Every integer $k\geq2$\\
Global identifiability (Theorem~\ref{thm:main}) & The above together with
\cite[Theorem~4.7, Corollary~5.7, Proposition~3.1]{OPT} & Every integer
$k\geq2$, generic parameters and initial values\\
\midrule
Coefficients, vector-field identity, recovery formulas &
Exact computer algebra, two separately written implementations
(Section~\ref{sec:checks}) & $k=1,\dots,12$ and $k=2,\dots,9$\\
Uniqueness for one trajectory & Exact power series and linear algebra &
$k=2,3,5$, one rational point\\
\midrule
Not claimed & Non-integer $k$; unknown $k$; an explicit description of
the exceptional set; practical identifiability from finite or noisy
data & \\
\bottomrule
\end{tabular}
\caption{Claims and their basis. The finite computations in the lower
part are checks of the general formulas and do not replace the proof.}
\label{tab:status}
\end{table}

The proof depends on three results of \cite{OPT}, whose hypotheses we
verify explicitly (Proposition~\ref{prop:charpres} and
Lemma~\ref{lem:noint}). Everything else is ordinary commutative algebra.

\section{Model and identifiability}
\label{sec:model}
Throughout, $k\geq2$ is a fixed integer, known to the observer. The
parameters are $\mu=(\alpha,\beta,\gamma)$, and we write
$K=\C(\alpha,\beta,\gamma)$ for the field of rational functions in them.
We write initial values as $S(0)$, $I(0)$, $R(0)$.

\paragraph{Reduction to two states.}
The equations for $S$ and $I$ in \eqref{eq:model} do not contain $R$.
If $(S,I,R)$ solves \eqref{eq:model} with $S+I+R=N_0$, then $(S,I)$ solves
\begin{equation}
 S'=-\frac{\beta SI}{u^k},\qquad
 I'=\frac{\beta SI}{u^k}-\gamma I,\qquad u=1+\alpha I,
 \qquad y=I,
 \label{eq:two}
\end{equation}
and conversely a solution of \eqref{eq:two} together with $R=N_0-S-I$
solves \eqref{eq:model}, since $(N_0-S-I)'=\gamma I$. Both have the same
output. Hence $\alpha,\beta,\gamma,S(0),I(0),R(0)$ are identifiable for
\eqref{eq:model} with known $N_0$ exactly when
$\alpha,\beta,\gamma,S(0),I(0)$ are identifiable for \eqref{eq:two}, and
then $R(0)=N_0-S(0)-I(0)$. The generic sets correspond under projection
to $(\mu,S(0),I(0))$. Pant, Saucedo and Pogudin use the same reduction
\cite[(S9.2)]{PSPsupp}, as does the authors' published code for $k=2$
\cite{PSPcode}. Only the last step uses $N_0$.

\paragraph{Identifiability.}
We use the definition of \cite[Definition~2.2]{OPT}, specialized to a
model without inputs. Write \eqref{eq:two} as $x'=f/Q$, $y=g/Q$ with
$x=(S,I)$, $Q=u^k$, $f=(-\beta SI,\ \beta SI-\gamma Iu^k)$ and $g=Iu^k$.
For $(\hat x_0,\hat\mu)\in\C^2\times\C^3$ with
$Q(\hat x_0,\hat\mu)\neq0$, let $Y(\hat x_0,\hat\mu)$ be the output of the
unique solution that is analytic near $t=0$.

\begin{definition}
\label{def:id}
A rational function $h(x_0,\mu)$ is \emph{identifiable} if there is a
nonempty Zariski open set $\Theta\subset\C^5$ such that, for every
$(\hat x_0,\hat\mu)\in\Theta$ at which $h$ and $Q^{-1}$ are defined, every
$(\tilde x_0,\tilde\mu)$ at which $h$ and $Q^{-1}$ are defined and with
$Y(\tilde x_0,\tilde\mu)=Y(\hat x_0,\hat\mu)$ satisfies
$h(\tilde x_0,\tilde\mu)=h(\hat x_0,\hat\mu)$. The model is
\emph{globally identifiable} if each parameter and each component of
$x_0$ is identifiable.
\end{definition}

Pant, Saucedo and Pogudin work over the reals with an open dense set of
admissible points \cite[Definition~1]{PSP}. As they note
\cite[Remark~4]{PSP}, identifiability over $\C$ implies identifiability
over $\mathbb R$: the real points of a nonempty Zariski open subset of
$\C^5$ form an open dense subset of $\mathbb R^5$, and a real competing
point is also a complex one. The set of real points of $\Theta$,
intersected with the region of positive parameters and populations, is
open and dense in that region.
Identifiability is a generic statement. For example, $I(0)=0$ gives the
zero output, and $S(0)=0$ gives pure exponential decay, from which
$\alpha$ and $\beta$ cannot be recovered.

\section{The input--output equation}
\label{sec:io}
Let $D$ be the derivation of $L=K(S,I)$ defined by \eqref{eq:two}, with
$D$ trivial on $K$; here $S$ and $I$ are independent indeterminates, and
$\alpha,\beta,\gamma$, being indeterminates, are nonzero in $K$.
This differential field is the field
$\operatorname{Frac}(\C(\mu)\{x,y\}/I_\Sigma)$ of \cite{OPT}: with
$I_\Sigma=[Qx'-f,\,Qy-g]:Q^\infty$ \cite[Notation~2.4]{OPT}, the ideal
$I_\Sigma$ is prime and consists of the differential polynomials that
vanish on all solutions \cite[Lemma~3.2]{HOPY}; see also the proof of
\cite[Corollary~5.7]{OPT}. It is the kernel of the differential
homomorphism $\C(\mu)\{x,y\}\to L$ sending $x$ to $(S,I)$ and $y$ to
$I$: every differential polynomial is congruent modulo $I_\Sigma$,
after multiplication by a power of $Q$, to a polynomial in $S,I$. If the
original element maps to zero in $L$, that polynomial is zero, because
$S$ and $I$ are algebraically independent over $K$; since $I_\Sigma$ is
saturated with respect to $Q$, the element lies in $I_\Sigma$.

Put $y=I$, $v=Dy$ and $w=D^2y$. The second equation of \eqref{eq:two}
and $D(S+I)=-\gamma I$ give
\begin{equation}
 S=\frac{(v+\gamma y)\,u^k}{\beta y},\qquad DS=-(v+\gamma y).
 \label{eq:S}
\end{equation}
Differentiating the first identity, substituting the second and
multiplying by $\beta y^2$, we obtain $P(y,v,w)=0$ on the model, where
\begin{equation}
 P=y u^k w-\bigl(1-(k-1)\alpha y\bigr)u^{k-1}v^2
 +\bigl(k\alpha\gamma y^2u^{k-1}+\beta y^2\bigr)v
 +\beta\gamma y^3 .
 \label{eq:P}
\end{equation}
Indeed, the derivative of $u^k$ contributes $k\alpha u^{k-1}v$; the
$v^2$ terms combine through $k\alpha y-u=-(1-(k-1)\alpha y)$, and the two
terms $\pm\gamma yu^kv$ cancel, leaving $k\alpha\gamma y^2u^{k-1}v$.

\begin{lemma}
\label{lem:irr}
The elements $y$ and $v$ of $L$ are algebraically independent over $K$,
and $P$ is irreducible in $K[y,v,w]$.
\end{lemma}

\begin{proof}
The map $(S,I)\mapsto(y,v)$ with $v=\beta Sy/u^k-\gamma y$ is inverted by
\eqref{eq:S}, so $K(S,I)=K(y,v)$, and $y,v$ are algebraically independent.
Write $P=Aw+B$ with $A=yu^k$ and $B\in K[y,v]$. The only irreducible
factors of $A$ in $K[y]$ are $y$ and $u=1+\alpha y$. From \eqref{eq:P},
\[
 B(0,v)=-v^2,\qquad
 B(-1/\alpha,v)=\frac{\beta}{\alpha^2}\Bigl(v-\frac{\gamma}{\alpha}\Bigr),
\]
where the second identity uses $k\geq2$, so that every term containing
$u^{k-1}$ vanishes at $u=0$. (For $k=1$ the value is also nonzero and the
lemma still holds; we restrict to $k\geq2$ only for brevity.) Both are nonzero, so neither $y$ nor $u$
divides $B$, and $\gcd(A,B)=1$. Thus $P$ is primitive and of degree one
in $w$ over $K[y,v]$; it is irreducible over $K(y,v)$ and hence, by Gauss's
lemma, in $K[y,v,w]$.
\end{proof}

\begin{proposition}
\label{prop:charpres}
Let $J=I_\Sigma\cap K\{y\}$, the kernel of the differential homomorphism
$K\{y\}\to L$, $y\mapsto I$. For the ranking $y<y'<y''<\cdots$, the set
$\{P\}$, with $y',y''$ in place of $v,w$, is a monic characteristic
presentation of $J$ in the sense of \cite[Definition~5.4]{OPT}.
\end{proposition}

\begin{proof}
$J$ is prime, being the kernel of a map to a field, and $P\in J$. In one
differential indeterminate an autoreduced set has at most one element.
The leader of $P$ is $y''$ and has degree one, so an element of $J$ of
lower rank would lie in $K[y,y']$; by Lemma~\ref{lem:irr} there is no
nonzero such element. Hence $\{P\}$ is a characteristic set of $J$
\cite[Definition~5.3]{OPT}. Its non-leading variables are $y,y'$. The
initial $A=yu^k$ lies in $K[y]\subset K[y,y']$, and, $P$ being irreducible
and of positive degree in $y''$, $P$ has no factor in $K[y,y']\setminus K$.
The coefficient of $yy''$ equals $1$, so $\{P\}$ is monic in the sense of
\cite[Definition~5.4]{OPT}.
\end{proof}

By \cite[Corollary~5.7]{OPT}, the field of input--output identifiable
functions is generated over $\C$ by the coefficients of $P$, as a
polynomial in $y,y',y''$.

\section{Coefficients and their inversion}
\label{sec:coef}
For a monomial $m$ in $y,v,w$ write $\coef{m}$ for its coefficient in
$P$, and set $\binom nr=0$ outside $0\leq r\leq n$.

\begin{proposition}
\label{prop:coef}
The coefficients of $P$, including those equal to zero, are
\begin{equation}
\begin{aligned}
 \coef{y^{j+1}w}&=\binom kj\alpha^j, && 0\leq j\leq k,\\
 \coef{y^{j+1}v}&=k\binom{k-1}{j-1}\gamma\alpha^j+\delta_{j1}\beta,
   && 1\leq j\leq k,\\
 \coef{y^{j}v^2}&=-\Bigl(\binom{k-1}{j}-(k-1)\binom{k-1}{j-1}\Bigr)\alpha^j,
   && 0\leq j\leq k,\\
 \coef{y^3}&=\beta\gamma,
\end{aligned}
\label{eq:table}
\end{equation}
and all other monomials have coefficient zero. In particular
\begin{equation}
 \coef{yw}=1,\quad \coef{y^2w}=k\alpha,\quad
 \coef{y^2v}=\beta+k\alpha\gamma,\quad
 \coef{y^3v}=k(k-1)\alpha^2\gamma,
 \label{eq:four}
\end{equation}
and, for $\alpha\gamma\neq0$,
\begin{equation}
 \alpha=\frac{\coef{y^2w}}{k},\qquad
 \gamma=\frac{\coef{y^3v}}{k(k-1)\alpha^2},\qquad
 \beta=\coef{y^2v}-k\alpha\gamma .
 \label{eq:inverse}
\end{equation}
\end{proposition}

\begin{proof}
Expand $u^k$ and $u^{k-1}$ in \eqref{eq:P} by the binomial theorem. The
$v^2$ coefficient is $-(u^{k-1}-(k-1)\alpha y\,u^{k-1})$, which gives the
third line; the second line comes from $k\alpha\gamma y^2u^{k-1}$ and
$\beta y^2$. Equations \eqref{eq:four} are the cases $j=0,1$ of the first
line and $j=1,2$ of the second, and \eqref{eq:inverse} solves them
successively.
\end{proof}

For instance, $\coef{yv^2}=0$ for every $k$. The denominators in
\eqref{eq:inverse} are $k$ and $k(k-1)\alpha^2$, which are nonzero in $K$
for every integer $k\geq2$. Since \eqref{eq:four} lies in the field
generated by the coefficients and $\alpha,\beta,\gamma$ are rational in
\eqref{eq:four}, Proposition~\ref{prop:charpres} and
\cite[Corollary~5.7]{OPT} show that $\alpha$, $\beta$ and $\gamma$ are
input--output identifiable, and that the field of input--output
identifiable functions is all of $K$.

For $k=1$ the coefficient $\coef{y^3v}$ vanishes, the field generated by
\eqref{eq:table} is $\C(\alpha,\beta+\alpha\gamma,\beta\gamma)$, and the
substitution $(\beta,\gamma)\mapsto(\alpha\gamma,\beta/\alpha)$ fixes all
coefficients. This agrees with \cite[Theorem~3.8]{PSP};
formulas~\eqref{eq:inverse} must not be applied for $k=1$.

\section{No rational first integrals}
\label{sec:noint}
Input--output identifiability does not by itself imply identifiability
\cite[Example~4.8]{OPT}. The bridge is the following lemma.

\begin{lemma}
\label{lem:noint}
Let $k\geq2$. The constants of the
differential field $(L,D)$, $L=K(S,I)$, are exactly $K$.
\end{lemma}

\begin{proof}
Use the coordinates $S$ and $u=1+\alpha I$, so that $L=K(S,u)$, and the
derivation $E=(u^k/I)\,D$. Since $u^k/I$ is a nonzero element of $L$,
$E$ and $D$ have the same constants. By \eqref{eq:two},
\begin{equation}
 ES=-\beta S,\qquad Eu=\alpha\beta S-\alpha\gamma u^k,
 \qquad
 E=-\beta S\,\partial_S+(\alpha\beta S-\alpha\gamma u^k)\,\partial_u .
 \label{eq:E}
\end{equation}
Let $0\neq F\in K(S,u)$ with $EF=0$. Embed $K(u)(S)$ into the field of
formal Laurent series $K(u)((S))$; no convergence is needed. Write
\[
 F=\sum_{n\geq m} f_n(u)\,S^n,\qquad f_m\neq0 .
\]
Because $E$ does not lower the $S$-order, it extends termwise to
$K(u)((S))$ compatibly with the embedding. The term $\alpha\beta
S\,\partial_u$ raises the order by one, so the coefficient of $S^m$ in
$EF$ is
\begin{equation}
 -\beta m f_m-\alpha\gamma u^k f_m'=0 .
 \label{eq:lead}
\end{equation}
Suppose $m\neq0$. Then $f_m'\neq0$, and
\[
 \frac{f_m'}{f_m}=-\frac{\beta m}{\alpha\gamma}\,u^{-k}.
\]
Write $f_m=u^r h$ with $r\in\mathbb Z$ and $h\in K(u)$ having $u$-adic
valuation zero. Then $f_m'/f_m=r/u+h'/h$ has $u$-adic valuation at least
$-1$, while the right-hand side has valuation $-k\leq-2$. This
contradiction shows $m=0$. Then \eqref{eq:lead} gives $f_0'=0$, so
$f_0\in K$ because the characteristic is zero. Now $F-f_0$ is again a
constant; if it were nonzero, its $S$-order would be a positive integer,
which the same argument excludes. Hence $F=f_0\in K$.
\end{proof}

The order of the pole, $k$, is the only place where the exponent enters.
For $k=1$ the right-hand side has a simple pole and the argument fails.
The lemma concerns the two-state model: with three states, $S+I+R$ is a
nonconstant rational first integral, and \cite[Theorem~4.7]{OPT} cannot
be applied directly.

\section{Proof of the main theorem}
\label{sec:proof}
\begin{theorem}
\label{thm:main}
Let $k\geq2$ be a fixed integer. For model \eqref{eq:model} with output
$y=I$ and known $N_0=S(0)+I(0)+R(0)$, the parameters $\alpha,\beta,\gamma$
and the initial values $S(0),I(0),R(0)$ are globally identifiable in the
sense of Definition~\ref{def:id}. The parameters and $S(0),I(0)$ are
identifiable without using $N_0$. Explicitly, $\alpha,\beta,\gamma$ are
given by \eqref{eq:inverse} from the normalized minimal input--output
equation, and
\begin{equation}
 S(t)=\frac{\bigl(y'(t)+\gamma y(t)\bigr)\bigl(1+\alpha y(t)\bigr)^k}
           {\beta y(t)},\qquad I(t)=y(t),\qquad R(t)=N_0-S(t)-I(t)
 \label{eq:states}
\end{equation}
near $t=0$, whenever $\beta I(0)(1+\alpha I(0))\neq0$.
\end{theorem}

\begin{proof}
By the reduction in Section~\ref{sec:model} it suffices to treat the
two-state model \eqref{eq:two}. By Lemma~\ref{lem:noint} the model has no
rational first integral, that is, the constants of
$\operatorname{Frac}(\C(\mu)\{x,y\}/I_\Sigma)=L$ are $\C(\mu)$. By
\cite[Theorem~4.7]{OPT}, a function of the parameters is identifiable if
and only if it is input--output identifiable. Section~\ref{sec:coef}
shows that $\alpha,\beta,\gamma$ are input--output identifiable; hence
they are identifiable.

For the initial values we use the criterion
\cite[Proposition~3.1]{OPT}: a rational function of $(x_0,\mu)$ is
identifiable if and only if its image in $L$, with $S(0),I(0)$ replaced
by $S,I$, lies in the subfield $M$ generated by the image of $\C\{y\}$,
that is, by $I,DI,D^2I,\dots$. Applied to $\alpha,\beta,\gamma$ the
criterion gives $\alpha,\beta,\gamma\in M$. By \eqref{eq:S},
$S=(DI+\gamma I)(1+\alpha I)^k/(\beta I)\in M$, and $I\in M$. Therefore
$S(0)$ and $I(0)$ are identifiable. Finally $R(0)=N_0-S(0)-I(0)$.

Formula \eqref{eq:states} is \eqref{eq:S} evaluated on a solution. The
condition $\beta I(0)(1+\alpha I(0))\neq0$ defines a nonempty Zariski open
set; on it $y$ and $1+\alpha y$ do not vanish near $t=0$, so the right-hand
side is defined there. Alternatively, analytically: once the parameters of any competing
point agree with the true ones on the generic set, the same output gives
the same right-hand side in \eqref{eq:states}, hence the same states.
\end{proof}

The theorem asserts uniqueness for generic, not for all, parameters and
initial values. We do not determine the smallest exceptional set. In
particular we do not claim that every point of an exceptional set is
non-identifiable.

\section{Finite exact checks}
\label{sec:checks}
The general argument above does not rely on computation. The following
exact computations were carried out as checks of the formulas; they are
finite and do not replace the proof.

\begin{enumerate}
\item \emph{Elimination from the vector field, $k=1,\dots,12$.} In a
Python~3 program using only the standard library (sparse polynomials
with integer coefficients and exact rationals), $S$ is eliminated from the relation that
is linear in $S$, $\beta Sy-(v+\gamma y)u^k=0$, and its total derivative
by a $2\times2$ resultant. The result agrees with \eqref{eq:table} in
every monomial. Substituting the vector field into $P$ and clearing
denominators gives the zero polynomial, as an identity rather than at
sample points. The four coefficients \eqref{eq:four}, the values $B(0,v)$
and $B(-1/\alpha,v)$ from Lemma~\ref{lem:irr}, and the coordinate form
\eqref{eq:E} of the derivation were also checked. For $k=1$ the
parameter points $(\alpha,\beta,\gamma)=(2,3,5)$ and $(2,10,3/2)$, related
by the substitution of Section~\ref{sec:coef}, give identical
coefficients; for $k=2,\dots,12$ they do not.
\item \emph{Comparison with published equations.} The input--output
equations printed in \cite[Section~S9]{PSPsupp} for $k=2,3,4$, obtained
there with \mbox{StructuralIdentifiability.jl}, coincide term by term with
$-P$.
\item \emph{A separately written implementation.} A second check, written
without reference to the first and using the computer algebra system
SymPy~1.13.3, verifies \eqref{eq:P} on the model for $k=2,3,5$ and for symbolic $k$
(treating $(1+\alpha I)^k$ as a separate symbol with the induced
derivative), the coefficient table \eqref{eq:table} and the inverse
\eqref{eq:inverse} for $k=2,\dots,9$, irreducibility of $P$ over
$\mathbb Q(\alpha,\beta,\gamma)$ for $k=2,3,5$, and the leading-order
identity \eqref{eq:lead}. For $k=2,3$ and random rational parameters it
finds no nonconstant polynomial first integral of degree at most $6$ in
$(S,u)$; this is consistent with Lemma~\ref{lem:noint} and is not a proof
of it.
\item \emph{One trajectory.} For $k=2,3,5$ and
$(\alpha,\beta,\gamma,S(0),I(0))=(1/3,2/5,1/7,5/2,1/2)$, the solution of
\eqref{eq:two} was computed as an exact rational power series to order
$40$. Among the polynomials $\sum c_{ijl}\,y^i(y')^j(y'')^l$ with
$i\leq k+1$ and $j+2l\leq2$, those vanishing on this trajectory form a
one-dimensional space; normalizing $c_{101}=1$ and applying
\eqref{eq:inverse} returns the original parameters exactly. Since the
input--output equation of any competing parameter point lies in the same
space of polynomials and has the same normalized coefficient, this
computation shows directly that this particular trajectory determines the
parameters.
\end{enumerate}
Scripts for these checks accompany this article, and the verification
record lists their outputs. Each check runs in a few seconds.

\section{Scope and remarks}
\label{sec:scope}
\paragraph{Integer exponents.}
Conjecture~2 of \cite{PSP} is stated for $k\geq2$, and the model text
calls $k\geq1$ the saturation exponent without restricting it to
integers. Theorem~\ref{thm:main} covers every fixed integer $k\geq2$. For
non-integer $k$ the right-hand side of \eqref{eq:model} is not rational,
the framework of \cite{OPT} does not apply as stated, and we make no claim.
We also do not treat $k$ as an unknown parameter to be recovered from the
data.

\paragraph{Relation to the large-$k$ sketch.}
Remark~1 of \cite[Section~S9]{PSPsupp} obtains global identifiability for
all $k$ outside an unspecified finite set from the software result for
the augmented model. Theorem~\ref{thm:main} shows that no integer $k\geq2$
is an exception to global identifiability, and
Proposition~\ref{prop:coef} gives the recovery formulas explicitly.

\paragraph{What is not addressed.}
Structural identifiability concerns exact, continuous observation. The
theorem says nothing about the stability of estimating $\alpha,\beta,\gamma$
from finitely many noisy measurements, about the fit of
\eqref{eq:model} to real epidemics, or about identifiability at the
non-generic parameter values excluded above.

\paragraph{Prior work.}
A literature search limited to the preprint versions of \cite{PSP}, the
authors' code and keyword searches found no other proof for general $k$;
it was not exhaustive, and we do not claim priority beyond it. The
argument is short once the results of \cite{OPT} are available, and it
may be known to specialists in a different form.

\section*{AI-use statement}
AI systems assisted with the derivation and checking of the proof, the
exact computer-algebra checks, independent readings of the proof and of
the manuscript, literature organization and manuscript drafting. The
author is responsible for the mathematical claims and their stated scope;
AI systems are not authors.

{\raggedright
\begin{thebibliography}{9}
\bibitem{PSP}
B. Pant, O. Saucedo and G. Pogudin,
``Uncovering identifiability of epidemiological models: basic
reproduction number and complementary data streams,''
medRxiv preprint, version 2 (posted 25 January 2026),
\href{https://doi.org/10.64898/2026.01.16.26344284}{doi:10.64898/2026.01.16.26344284},
\url{https://www.medrxiv.org/content/10.64898/2026.01.16.26344284v2}.

\bibitem{PSPsupp}
B. Pant, O. Saucedo and G. Pogudin,
Supplementary material to \cite{PSP}, version 2, Sections S8 and S9,
\url{https://www.medrxiv.org/content/medrxiv/early/2026/01/25/2026.01.16.26344284/DC1/embed/media-1.pdf}.

\bibitem{PSPcode}
B. Pant, code accompanying \cite{PSP}, GitHub repository
\url{https://github.com/Binod-Pant/structural-identifiability-complementary},
file \texttt{behavior~model/behavior\_local\_power\_two\_julia.jl},
commit \nolinkurl{c2a43c695d4055b97b1d9dd3cedd6e7e56b7b8ff}.

\bibitem{OPT}
A. Ovchinnikov, G. Pogudin and P. Thompson,
``Parameter identifiability and input--output equations,''
\emph{Appl. Algebra Engrg. Comm. Comput.} \textbf{34}(2), 165--182 (2023;
published online 2021),
\href{https://doi.org/10.1007/s00200-021-00486-8}{doi:10.1007/s00200-021-00486-8};
\href{https://arxiv.org/abs/2007.14787v2}{arXiv:2007.14787v2} (2020).
Theorem, definition and corollary numbers in this article refer to
arXiv:2007.14787v2.

\bibitem{HOPY}
H. Hong, A. Ovchinnikov, G. Pogudin and C. Yap,
``Global identifiability of differential models,''
\emph{Comm. Pure Appl. Math.} \textbf{73}(9), 1831--1879 (2020),
\href{https://doi.org/10.1002/cpa.21921}{doi:10.1002/cpa.21921};
\href{https://arxiv.org/abs/1801.08112v5}{arXiv:1801.08112v5}.
The lemma number refers to arXiv:1801.08112v5.
\end{thebibliography}
}
\end{document}
